\section{Introduction}\label{sec:introduction}
Let \(V\) be a finite-dimensional real inner-product space with an
orthogonal representation \(\rho:O(3)\to O(V)\). Let
\(\mathcal X=\mathcal S'(\R^3;V)\), equipped with the sigma-algebra
generated by evaluations. For \(g(x)=Qx+b\), use
\begin{equation}\label{eq:euclidean-action}
 (g\cdot f)(x)=\rho(Q)f(Q^{-1}(x-b)),\qquad
 (g\cdot F)(f)=F(g^{-1}\cdot f),\qquad
 \widehat g A(F)=A(g^{-1}\cdot F).
\end{equation}
Write \(U(a)=\widehat{\tau_a}\).
For open \(O\subset\R^3\), let \(\Sigma_O\) be generated by
\(F(f)\) with real \(f\in C_c^\infty(O;V)\).
For a reflection \(r\), put \(\mathcal R_r=\widehat r\).
An \(r\)-invariant law \(\lambda\) is reflection positive on an
open side \(H^+\) if
\[
 \E_\lambda[\overline{\mathcal R_r A}\,A]\ge0
\]
for finite character sums \(A=\sum_bz_be^{iF(f_b)}\) with
real \(f_b\in C_c^\infty(H^+;V)\).
The positive sides below are chosen by making the displayed linear
coordinate function positive.

\begin{theorem}\label{thm:components-clustering}
Let \(\lambda\) be a translation-invariant probability law on
\(\mathcal X\), invariant and reflection positive at the nine planes
\begin{equation}\label{eq:nine-planes}
 x_j=0\quad(1\le j\le3),\qquad
 x_j-x_k=0,\quad x_j+x_k=0\quad(1\le j<k\le3).
\end{equation}
Let \(\mathcal I\) be the translation-invariant sigma-algebra,
with invariance understood modulo \(\lambda\)-null sets, and let
\(P=\E_\lambda[\,\cdot\mid\mathcal I]\). Then
\begin{equation}\label{eq:main-relative-clustering}
 \langle A,U(a)B\rangle_{L^2(\lambda)}
 \longrightarrow\langle A,PB\rangle_{L^2(\lambda)}
 \qquad(|a|\longrightarrow\infty)
\end{equation}
for every \(A,B\in L^2(\lambda)\).

There is a unique translation-ergodic decomposition
\(\lambda=\int\lambda_\xi\,\dd\nu(\xi)\), understood as the
distribution of its conditional law on \(\mathcal X\).
On one common \(\nu\)-full set, every \(\lambda_\xi\) is
translation ergodic, invariant and reflection positive at every
real translate of the planes in \eqref{eq:nine-planes}, and
\begin{equation}\label{eq:main-component-clustering}
 \langle A,U(a)B\rangle_{L^2(\lambda_\xi)}
 \longrightarrow
 \overline{\E_{\lambda_\xi}A}\,\E_{\lambda_\xi}B
 \qquad(|a|\longrightarrow\infty)
\end{equation}
for every \(A,B\in L^2(\lambda_\xi)\).
Reflection positivity holds on the entire positive-side
\(L^2(\lambda_\xi)\) space. These component laws are cubically invariant.

If in addition, for some fixed \(C\ge0\),
\begin{equation}\label{eq:probability-G}
 \E_\lambda e^{F(f)}\le e^{C\|f\|_2^2/2}
 \qquad(f\in\mathcal S(\R^3;V)\text{ real}),
\end{equation}
then the common full-measure set above may be chosen so that every
component in it satisfies \eqref{eq:probability-G} with the same
\(C\), simultaneously for all real tests.
Any specified measurable event of full \(\lambda\)-measure has
full measure in almost every component.
\end{theorem}

The proof follows the reflection and component arguments of
\cite[Section~7]{QED3Manuscript}. They depend on the law and its
geometric action, so the proof applies to the finite component
space \(V\) above. No field equation or Gaussian estimate is used
for \eqref{eq:main-relative-clustering} or the reflection properties
of the components.

The spectral step is essential for the all-direction limit.
Coordinate reflection positivity constructs a nonnegative generator
\(H_j\) and commuting transverse momenta \(\mathsf P_i\).
A second reconstruction at a diagonal mirror gives a bound on
Laplace transforms of their joint spectral measure. We use
transforms whose suprema on the positive real and imaginary axes
have bounded product, while their modulus grows exponentially
outside the sector \(E\ge|p_i|\).
This proves \(H_j\ge|\mathsf P_i|\).
At zero energy all transverse momenta therefore vanish.
It follows that invariance under one coordinate translation group
is already invariance under all translations. Semigroup convergence
then gives \eqref{eq:main-relative-clustering}, uniformly in
transverse displacements.

For the component assertions, a bounded translation-invariant
observable has a representative in either half-space and is fixed
by its reflection. Applying reflection positivity to \(\sqrt h A\),
with \(h\) a nonnegative invariant weight, yields a conditional
reflection inequality. Countable determining character algebras
put the resulting properties on one set of component laws.
The optional Gaussian assertion uses a separate estimate:
spatial averages of \(e^{F(f)}\) have moments controlled by
translated \(L^2\) overlaps of \(f\). Those overlaps vanish at
large separation, which preserves the original constant \(C\)
after conditioning.

The conclusion is qualitative clustering. It asserts no positive
spectral gap. The hypotheses also do not assert invariance under
arbitrary rotations or field-sign change in each component.
